\chapter{Arhitectura Sistemului Hardware}
\label{chap:arhitectura_hardware}

\section{Topologia Rețelei și Configurația Bancului de Test}
Sistemul dezvoltat utilizează o topologie distribuită formată din trei noduri inteligente bazate pe SoC-ul (System-on-Chip) Espressif ESP32 și un analizor radio programabil Ubertooth One în banda ISM de 2.4 GHz.

Cele trei noduri ESP32 sunt poziționate într-o geometrie triunghiulară în planul încăperii de test ($5\text{ m} \times 5\text{ m}$), formând trei legături radio bidirecționale distincte ($L_{12}$, $L_{23}$, $L_{31}$):
\begin{itemize}
    \item \textbf{Nodul 1 (Receptor / Colector principal):} Modul ESP32-WROOM-32 (adresa MAC \texttt{fc:e8:c0:e1:20:bc}) conectat prin interfață USB-UART la stația de procesare pe portul \texttt{/dev/ttyUSB1}. Poziție carteziană: $(-1.5\text{ m}, +1.0\text{ m})$.
    \item \textbf{Nodul 2 (Receptor secundar):} Modul ESP32-WROOM-32 (adresa MAC \texttt{fc:e8:c0:e1:1e:d0}) conectat pe portul \texttt{/dev/ttyUSB0}. Poziție carteziană: $(+1.5\text{ m}, +1.0\text{ m})$.
    \item \textbf{Nodul 3 (Transmițător beacon autonom):} Modul ESP32 (adresa MAC \texttt{fc:e8:c0:e1:1a:ec}) alimentat direct dintr-un adaptor de priză la rețea ($5\text{ V}$, $2\text{ A}$). Poziție carteziană: $(0.0\text{ m}, -1.5\text{ m})$. Emite cadre de sondare (Probe Request) la o frecvență ridicată de $250\text{ Hz}$ în regim continuu.
    \item \textbf{Ubertooth One:} Dispozitiv SDR de laborator (identificator USB \texttt{1d50:6002}) utilizat pentru supravegherea spectrului radio în 2.4 GHz, monitorizarea perturbațiilor provocate de dispozitivele Bluetooth concurente și injectarea de cadre de calibrare pilot.
\end{itemize}

\section{Firmware-ul Dedicat ESP-IDF v5 pentru Achiziție CSI}
Pe plăcile ESP32 a fost dezvoltat și încărcat un firmware custom compilat cu instrumentul oficial Espressif IoT Development Framework (ESP-IDF v5.1).

\subsection{Configurarea Modului Promiscuu și a Modulului CSI}
Pentru a permite extragerea matricei CSI la nivel de nivel fizic (PHY), a fost necesară configurarea fișierului \texttt{sdkconfig} cu opțiunea \texttt{CONFIG\_ESP\_WIFI\_CSI\_ENABLED=y}. În caz contrar, apelul funcției \texttt{esp\_wifi\_set\_csi\_config()} returnează codul de eroare critic \texttt{0xffffffff (ESP\_FAIL)}.

Firmware-ul inițializează subsistemul Wi-Fi în mod promiscuu complet (\texttt{WIFI\_PROMIS\_FILTER\_MASK\_ALL}):
\begin{verbatim}
wifi_promiscuous_filter_t filter = {
    .filter_mask = WIFI_PROMIS_FILTER_MASK_ALL
};
esp_wifi_set_promiscuous_filter(&filter);
esp_wifi_set_promiscuous(true);

wifi_csi_config_t csi_config = {
    .lltf_en = true,
    .htltf_en = true,
    .stbc_htltf2_en = true,
    .ltf_merge_en = true,
    .channel_filter_en = true,
    .manu_scale = false,
    .shift = false
};
esp_wifi_set_csi_config(&csi_config);
esp_wifi_set_csi_rx_cb(&wifi_csi_rx_callback, NULL);
\end{verbatim}

\subsection{Optimizarea Sarcinii FreeRTOS și Cadența de Emisie}
Pe nodul transmițător (Nod 3), rutina de generare a pachetelor rulează într-o sarcină dedicată FreeRTOS pe nucleul de execuție Core 1. Inițial, apelul \texttt{pdMS\_TO\_TICKS(4)} evalua la valoarea 0 când ceasul de sistem este configurat la \texttt{CONFIG\_FREERTOS\_HZ=100}, generând o panică de kernel de tipul \texttt{assert failed: xTaskDelayUntil}. Problema a fost soluționată prin apelarea directă a funcției \texttt{vTaskDelay(1)} combinată cu temporizatorul hardware de înaltă precizie \texttt{esp\_timer\_get\_time()}, asigurând o cadență constantă și stabilă de emisie.

\section{Formatul Pachetelor Seriale UART (138 Octeți)}
Pentru a asigura o lățime de bandă suficientă pentru eșantionarea a 250 de pachete pe secundă fără întârzieri de buffer, porturile seriale USB sunt configurate la rata de transfer de \textbf{921600 baud}.

Fiecare pachet binar emis pe magistrala serială UART are o lungime fixă de exact 138 de octeți, structurat conform specificației din Tabelul \ref{tab:packet_structure}.

\begin{table}[H]
\centering
\begin{tabular}{@{}llll@{}}
\toprule
\textbf{Offset [octeți]} & \textbf{Câmp} & \textbf{Tip Dată} & \textbf{Descriere Funcțională} \\ \midrule
0 -- 1 & \texttt{Magic Header} & \texttt{uint16\_t} & Număr magic de sincronizare: \texttt{0xA55A} \\
2 & \texttt{Node ID} & \texttt{uint8\_t} & Identificator numeric al nodului emițător (1..3) \\
3 & \texttt{RSSI} & \texttt{int8\_t} & Nivel putere semnal recepționat (dBm) \\
4 -- 7 & \texttt{Timestamp} & \texttt{uint32\_t} & Marcaj temporal hardware de la boot ($\mu$s) \\
8 -- 9 & \texttt{Sequence No} & \texttt{uint16\_t} & Număr secvențial pachet (detecție pierderi) \\
10 -- 137 & \texttt{CSI Raw Data} & \texttt{int8\_t[128]} & 64 perechi $[I_k, Q_k]$ pentru subpurtătoare \\ \bottomrule
\end{tabular}
\caption{Structura pachetului binar de telemetrie CSI pe magistrala serială UART.}
\label{tab:packet_structure}
\end{table}

\subsection{Decodarea și Extragerea celor 52 de Subpurtătoare Active}
În blocul de 128 de octeți al câmpului de date CSI, receptorul raportează 64 de valori complexe ordonate după schema:
\begin{equation}
    \text{IQ\_Data} = [I_0, Q_0, I_1, Q_1, \dots, I_{63}, Q_{63}]
\end{equation}
Pentru standardul IEEE 802.11 HT20, subpurtătoarele marginale nule și frecvența purtătoare continuă DC (indicele 32) sunt eliminate conform standardului. Vectorul final de 52 de coeficienți activi este obținut prin concatenarea feliilor:
\begin{equation}
    CSI_{52} = [CSI[6:32], CSI[33:59]]
\end{equation}
Amplitudinea scalară a fiecărei subpurtătoare $k \in \{0, \dots, 51\}$ se calculează conform relației:
\begin{equation}
    |H_k| = \sqrt{I_k^2 + Q_k^2}
\end{equation}
Această matrice bidimensională de dimensiune $N_{eșantioane} \times 52$ intră direct în motorul de procesare digitală a semnalului.
